% GENERATED FILE. Edit canonical lessons, then run npm run generate:books.
% canonical-source-hash: fnv1a64:1ac408af26fbf793
% canonical-lessons: HI-C552-jabki, HI-C552-bhalehi, HI-C552-baharhal, HI-C552-khair, HI-C552-magar, HI-R552-review-546-552

\chapter{Even So}
\label{ch:hi-even-so}
\hlchaptermodality{552}

\begin{chapteropening}
\textbf{By the end of this chapter:} \emph{I can link ideas with whereas, even though, anyway, well and but.}

Complete the last lesson of chapter 552: I can link ideas with whereas, even though, anyway, well and but.
\end{chapteropening}

\section[\texorpdfstring{\hi{जबकि} (jabki)}{जबकि (jabki)}]{\hi{जबकि} (jabki) — whereas; while}

\label{lesson:HI-C552-jabki}



\begin{quote}
Before the new one: say the Hindi for stubborn, then the Hindi for force, coercion.
\end{quote}

\subsection*{You'll want to know: \hi{जबकि}}
\textbf{\hi{जबकि}} — \emph{jabki} — \textquotedblleft{}whereas; while\textquotedblright{}.

\begin{quote}
\hi{मुझे} \hi{चाय} \hi{पसंद} \hi{है}, \hi{जबकि} \hi{मीरा} \hi{को} \hi{दूध} \hi{पसंद} \hi{है।}
\end{quote}

\begin{quote}
\emph{Mujhe cāy pasand hai, jabki Mīrā ko dūdh pasand hai.}
\end{quote}

\begin{quote}
I like tea, whereas Mira likes milk.
\end{quote}

\begin{cousinweb}
\textbf{\hi{जबकि}} (\emph{jabki}) joins \emph{jab}, \textquotedblleft{}when\textquotedblright{}, and \emph{ki}, \textquotedblleft{}that\textquotedblright{}, which is usually traced to Persian \emph{ki}.
\end{cousinweb}

\begin{grammarlens}[title={What you've built}]
One more word for saying what happens if.
\end{grammarlens}

\subsection*{Your turn}
\begin{itemize}
\raggedright
  \item \emph{Say it:} \emph{jabki}
  \item \emph{Say it:} \emph{jabki}, once more
  \item \emph{Say it:} \emph{Mujhe cāy pasand hai, jabki Mīrā ko dūdh pasand hai.} — I like tea, whereas Mira likes milk.
  \item \emph{Recall:} say the Hindi for stubborn, then the Hindi for force, coercion, then say \emph{jabki} again
\end{itemize}

\begin{tcolorbox}[breakable,colback=teal!4,colframe=teal!35!black,title={Before you move on}]
What does \hi{जबकि} mean? (\textquotedblleft{}Whereas; while\textquotedblright{}.) Say it once more.
\end{tcolorbox}
\section[\texorpdfstring{\hi{भले} \hi{ही} (bhale hī)}{भले ही (bhale hī)}]{\hi{भले} \hi{ही} (bhale hī) — even though; even if}

\label{lesson:HI-C552-bhalehi}



\begin{quote}
Before the new one: say the Hindi for force, coercion, then the Hindi for whereas.
\end{quote}

\subsection*{You'll want to know: \hi{भले} \hi{ही}}
\textbf{\hi{भले} \hi{ही}} — \emph{bhale hī} — \textquotedblleft{}even though; even if\textquotedblright{}.

\begin{quote}
\hi{भले} \hi{ही} \hi{रास्ता} \hi{लंबा} \hi{है}, \hi{फिर} \hi{भी} \hi{हम} \hi{चलकर} \hi{जाएँगे।}
\end{quote}

\begin{quote}
\emph{Bhale hī rāstā lambā hai, phir bhī ham calkar jāeṅge.}
\end{quote}

\begin{quote}
Even though the way is long, we will go on foot.
\end{quote}

\begin{cousinweb}
\textbf{\hi{भले} \hi{ही}} (\emph{bhale hī}) is \emph{bhale}, a form of \emph{bhalā}, \textquotedblleft{}good\textquotedblright{}, which CDIAL derives from Sanskrit \emph{bhalla}, \textquotedblleft{}auspicious\textquotedblright{} (CDIAL 9408), with the emphatic \emph{hī}: \textquotedblleft{}well and good, even so\textquotedblright{}.
\end{cousinweb}

\begin{grammarlens}[title={What you've built}]
One more word for saying what happens if.
\end{grammarlens}

\subsection*{Your turn}
\begin{itemize}
\raggedright
  \item \emph{Say it:} \emph{bhale hī}
  \item \emph{Say it:} \emph{bhale hī}, once more
  \item \emph{Say it:} \emph{Bhale hī rāstā lambā hai, phir bhī ham calkar jāeṅge.} — Even though the way is long, we will go on foot.
  \item \emph{Recall:} say the Hindi for force, coercion, then the Hindi for whereas, then say \emph{bhale hī} again
\end{itemize}

\begin{tcolorbox}[breakable,colback=teal!4,colframe=teal!35!black,title={Before you move on}]
What does \hi{भले} \hi{ही} mean? (\textquotedblleft{}Even though; even if\textquotedblright{}.) Say it once more.
\end{tcolorbox}
\section[\texorpdfstring{\hi{बहरहाल} (baharhāl)}{बहरहाल (baharhāl)}]{\hi{बहरहाल} (baharhāl) — anyway; in any case}

\label{lesson:HI-C552-baharhal}



\begin{quote}
Before the new one: say the Hindi for whereas, then the Hindi for even though.
\end{quote}

\subsection*{You'll want to know: \hi{बहरहाल}}
\textbf{\hi{बहरहाल}} — \emph{baharhāl} — \textquotedblleft{}anyway; in any case\textquotedblright{}.

\begin{quote}
\hi{रास्ता} \hi{लंबा} \hi{था}, \hi{बहरहाल} \hi{हम} \hi{समय} \hi{पर} \hi{पहुँच} \hi{गए।}
\end{quote}

\begin{quote}
\emph{Rāstā lambā thā, baharhāl ham samay par pahũc gae.}
\end{quote}

\begin{quote}
The way was long; anyway, we arrived on time.
\end{quote}

\begin{cousinweb}
\textbf{\hi{बहरहाल}} (\emph{baharhāl}) is usually traced to Persian \emph{ba-har-hāl}, \textquotedblleft{}in every state\textquotedblright{}.
\end{cousinweb}

\begin{grammarlens}[title={What you've built}]
One more word for saying what happens if.
\end{grammarlens}

\subsection*{Your turn}
\begin{itemize}
\raggedright
  \item \emph{Say it:} \emph{baharhāl}
  \item \emph{Say it:} \emph{baharhāl}, once more
  \item \emph{Say it:} \emph{Rāstā lambā thā, baharhāl ham samay par pahũc gae.} — The way was long; anyway, we arrived on time.
  \item \emph{Recall:} say the Hindi for whereas, then the Hindi for even though, then say \emph{baharhāl} again
\end{itemize}

\begin{tcolorbox}[breakable,colback=teal!4,colframe=teal!35!black,title={Before you move on}]
What does \hi{बहरहाल} mean? (\textquotedblleft{}Anyway; in any case\textquotedblright{}.) Say it once more.
\end{tcolorbox}
\section[\texorpdfstring{\hi{ख़ैर} (khair)}{ख़ैर (khair)}]{\hi{ख़ैर} (khair) — well; anyway (moving on)}

\label{lesson:HI-C552-khair}



\begin{quote}
Before the new one: say the Hindi for even though, then the Hindi for anyway.
\end{quote}

\subsection*{You'll want to know: \hi{ख़ैर}}
\textbf{\hi{ख़ैर}} — \emph{khair} — \textquotedblleft{}well; anyway (moving on)\textquotedblright{}.

\begin{quote}
\hi{ख़ैर}, \hi{अब} \hi{चलिए}, \hi{देर} \hi{हो} \hi{रही} \hi{है।}
\end{quote}

\begin{quote}
\emph{Khair, ab calie, der ho rahī hai.}
\end{quote}

\begin{quote}
Well, let's go now; it's getting late.
\end{quote}

\begin{cousinweb}
\textbf{\hi{ख़ैर}} (\emph{khair}) is usually traced to Arabic \emph{khayr}, \textquotedblleft{}good, well-being\textquotedblright{}, through Persian.
\end{cousinweb}

\begin{grammarlens}[title={What you've built}]
One more word for saying what happens if.
\end{grammarlens}

\subsection*{Your turn}
\begin{itemize}
\raggedright
  \item \emph{Say it:} \emph{khair}
  \item \emph{Say it:} \emph{khair}, once more
  \item \emph{Say it:} \emph{Khair, ab calie, der ho rahī hai.} — Well, let's go now; it's getting late.
  \item \emph{Recall:} say the Hindi for even though, then the Hindi for anyway, then say \emph{khair} again
\end{itemize}

\begin{tcolorbox}[breakable,colback=teal!4,colframe=teal!35!black,title={Before you move on}]
What does \hi{ख़ैर} mean? (\textquotedblleft{}Well; anyway (moving on)\textquotedblright{}.) Say it once more.
\end{tcolorbox}
\section[\texorpdfstring{\hi{मगर} (magar)}{मगर (magar)}]{\hi{मगर} (magar) — but (in a lighter register); yet}

\label{lesson:HI-C552-magar}



\begin{quote}
Before the new one: say the Hindi for anyway, then the Hindi for well.
\end{quote}

\subsection*{You'll want to know: \hi{मगर}}
\textbf{\hi{मगर}} — \emph{magar} — \textquotedblleft{}but (in a lighter register); yet\textquotedblright{}.

\begin{quote}
\hi{मैं} \hi{बहुत} \hi{थका} \hi{था}, \hi{मगर} \hi{मैंने} \hi{खाना} \hi{बनाया।}
\end{quote}

\begin{quote}
\emph{Maiṁ bahut thakā thā, magar maiṁne khānā banāyā.}
\end{quote}

\begin{quote}
I was very tired, but I cooked the meal.
\end{quote}

Means the same as \hi{लेकिन}; common in speech and in songs.

\begin{cousinweb}
\textbf{\hi{मगर}} (\emph{magar}) is usually traced to Persian \emph{magar}, \textquotedblleft{}but, unless\textquotedblright{}.
\end{cousinweb}

\begin{grammarlens}[title={What you've built}]
One more word for saying what happens if.
\end{grammarlens}

\subsection*{Your turn}
\begin{itemize}
\raggedright
  \item \emph{Say it:} \emph{magar}
  \item \emph{Say it:} \emph{magar}, once more
  \item \emph{Say it:} \emph{Maiṁ bahut thakā thā, magar maiṁne khānā banāyā.} — I was very tired, but I cooked the meal.
  \item \emph{Recall:} say the Hindi for anyway, then the Hindi for well, then say \emph{magar} again
\end{itemize}

\begin{tcolorbox}[breakable,colback=teal!4,colframe=teal!35!black,title={Before you move on}]
What does \hi{मगर} mean? (\textquotedblleft{}But (in a lighter register); yet\textquotedblright{}.) Say it once more.
\end{tcolorbox}
\section[(dialogue)]{Looking back: from Fate and Chance to Even So}

\label{lesson:HI-R552-review-546-552}



\begin{quote}
Say the Hindi for an omen (a good sign) and signs.
\end{quote}

\subsection*{Your turn}
Say these aloud:

\begin{itemize}
\raggedright
  \item an omen (a good sign), signs, unlucky, good fortune, chance
  \item alert, watchful, safekeeping, watch, an inkling, an instruction
  \item a period, a term, the end, laziness, to squander
  \item an option, a dilemma, indecision, preference, to hesitate, to hold back
  \item one's own wish, of one's own choosing, a ban, a prohibition, favourite, a selection
  \item approved, free, helpless, stubborn, force, coercion
  \item whereas, even though, anyway, well, but (in a lighter register)
\end{itemize}

\begin{tcolorbox}[breakable,colback=teal!4,colframe=teal!35!black,title={Before you move on}]
An omen (a good sign)? (\textbf{\hi{शगुन}}, \emph{śagun}.) Signs? (\textbf{\hi{आसार}}, \emph{āsār}.) Unlucky? (\textbf{\hi{मनहूस}}, \emph{manhūs}.) Good fortune? (\textbf{\hi{सौभाग्य}}, \emph{saubhāgya}.) Chance? (\textbf{\hi{इत्तफ़ाक़}}, \emph{ittafāq}.) Alert, watchful? (\textbf{\hi{चौकन्ना}}, \emph{caukannā}.) Safekeeping? (\textbf{\hi{हिफ़ाज़त}}, \emph{hifāzat}.) Watch? (\textbf{\hi{निगरानी}}, \emph{nigrānī}.) An inkling? (\textbf{\hi{भनक}}, \emph{bhanak}.) An instruction? (\textbf{\hi{हिदायत}}, \emph{hidāyat}.) A period? (\textbf{\hi{अवधि}}, \emph{avadhi}.) A term? (\textbf{\hi{मियाद}}, \emph{miyād}.) The end? (\textbf{\hi{अंत}}, \emph{ant}.) Laziness? (\textbf{\hi{आलस}}, \emph{ālas}.) To squander? (\textbf{\hi{लुटाना}}, \emph{luṭānā}.) An option? (\textbf{\hi{विकल्प}}, \emph{vikalp}.) A dilemma? (\textbf{\hi{दुविधा}}, \emph{duvidhā}.) Indecision? (\textbf{\hi{असमंजस}}, \emph{asmañjas}.) Preference? (\textbf{\hi{तरजीह}}, \emph{tarjīh}.) To hesitate, to hold back? (\textbf{\hi{हिचकना}}, \emph{hicaknā}.) One's own wish? (\textbf{\hi{मर्ज़ी}}, \emph{marzī}.) Of one's own choosing? (\textbf{\hi{मनपसंद}}, \emph{manpasand}.) A ban, a prohibition? (\textbf{\hi{मनाही}}, \emph{manāhī}.) Favourite? (\textbf{\hi{पसंदीदा}}, \emph{pasandīdā}.) A selection? (\textbf{\hi{चयन}}, \emph{cayan}.) Approved? (\textbf{\hi{मंज़ूर}}, \emph{manzūr}.) Free? (\textbf{\hi{आज़ाद}}, \emph{āzād}.) Helpless? (\textbf{\hi{लाचार}}, \emph{lācār}.) Stubborn? (\textbf{\hi{अड़ियल}}, \emph{aṛiyal}.) Force, coercion? (\textbf{\hi{ज़बरदस्ती}}, \emph{zabardastī}.) Whereas? (\textbf{\hi{जबकि}}, \emph{jabki}.) Even though? (\textbf{\hi{भले} \hi{ही}}, \emph{bhale hī}.) Anyway? (\textbf{\hi{बहरहाल}}, \emph{baharhāl}.) Well? (\textbf{\hi{ख़ैर}}, \emph{khair}.) But (in a lighter register)? (\textbf{\hi{मगर}}, \emph{magar}.)
\end{tcolorbox}
